\documentclass[11pt]{article}
\usepackage[margin=1in]{geometry}
\usepackage{amsmath,amssymb,amsthm,mathtools}
\usepackage{hyperref}

\theoremstyle{plain}
\newtheorem{proposition}{Proposition}
\newtheorem{lemma}{Lemma}
\newtheorem{corollary}{Corollary}
\newtheorem{conjecture}{Conjecture}
\theoremstyle{definition}
\newtheorem{definition}{Definition}
\theoremstyle{remark}
\newtheorem{remark}{Remark}

\newcommand{\C}{\mathbb{C}}
\newcommand{\R}{\mathbb{R}}
\newcommand{\E}{\mathbb{E}}
\newcommand{\dt}{\partial_t}
\newcommand{\dx}{\partial_x}
\newcommand{\abs}[1]{\lvert #1\rvert}
\newcommand{\LP}{\mathcal{LP}}

\title{A Gershgorin scaling law for quartic decoherence of zeros}
\author{}
\date{\today}

\begin{document}
\maketitle

\begin{abstract}
Consider the fourth-order heat flow $\dt H=\dx^4H$ acting on a real function with
only real zeros. There are two natural criteria for when the flow lifts a close
real pair of zeros off the axis: a \emph{floor} criterion (the pair drowns when
$tH_0^{(4)}$ overwhelms the small extremum the pair sits over) and a
\emph{drift} criterion (writing the flow as a Gaussian average of second-order
heat flows over signed times, the pair decoheres when the displacement
$\abs v\sqrt{2t}$ across the averaging window reaches the half-gap, giving
$t_c=h^2/2v^2$). We show these are one law in Gershgorin form. For an isolated
close pair at half-gap $h$ embedded in a field of remaining zeros, the lift-off
time is exactly
\[
   t_{\mathrm{drown}}=\frac{h^2}{24\,e_2}\,(1+O(h^2)),
\]
where $e_2$ is the second elementary symmetric function of the inverse distances
from the collision site to the other zeros. Lift-off requires $e_2>0$ (net
one-sidedness/drift), which makes a symmetric lattice stable and an edge
configuration unstable; in the drift-dominated regime $e_2\approx p_1^2/2$ this
reduces to the drift form $t_c\asymp h^2/p_1^2$. The common law
$t_{\mathrm{drown}}\asymp h^2/(\text{crowding})$ makes gap compression explicit:
where local gaps close while the surrounding crowding stays bounded below,
lift-off times shrink to zero. We instantiate this at a random-matrix spectral
edge, where the one-sided field forces instability with edge gaps $h\asymp
N^{-2/3}$. We make no claims about the Riemann hypothesis.
\end{abstract}

\section{Setup and the two criteria}
Let $H_0$ be real entire (or a real polynomial) with only simple real zeros, and
evolve it by the fourth-order heat flow
\begin{equation}\label{eq:flow}
  \dt H=\dx^{4}H,\qquad H|_{t=0}=H_0,
\end{equation}
so that locally $H_t=H_0+t\,H_0^{(4)}+\tfrac{t^2}{2}H_0^{(8)}+\cdots$ (a finite
sum for polynomials). Two criteria govern when a real pair of zeros becomes a
complex-conjugate pair.

\paragraph{Floor criterion.} If two real zeros bound an extremum of $H_0$ at
$x_*$ (so $H_0'(x_*)=0$), the value $H_t(x_*)$ moves linearly at first order; the
pair merges and lifts off when $H_t(x_*)$ changes sign,
\begin{equation}\label{eq:floor}
  t_{\mathrm{drown}}=-\frac{H_0(x_*)}{H_0^{(4)}(x_*)} .
\end{equation}

\paragraph{Drift criterion.} The fourth-order flow admits the Gaussian-average
representation
\begin{equation}\label{eq:avg}
  H_t=\E_{s\sim N(0,2t)}\bigl[\,e^{-s\dx^2}H_0\,\bigr],
\end{equation}
since $\E_{s\sim N(0,2t)}e^{su^2}=e^{tu^4}$ on the Fourier side: the quartic flow
is the second-order heat flow run for a random, two-sided Gaussian time $s$. A
zero $x_j$ of the second-order flow $e^{-s\dx^2}H_0$ moves with velocity
$v=H_0''/H_0'$ (evaluated at $x_j$); across the averaging window of width
$\sqrt{2t}$ its displacement is $\abs v\sqrt{2t}$, and two zeros at half-gap $h$
decohere into the average when this reaches $h$, i.e.\ at
\begin{equation}\label{eq:drift}
  t_c=\frac{h^2}{2v^2}.
\end{equation}
Both \eqref{eq:floor} and \eqref{eq:drift} scale as $h^2$. We now show they are
the same Gershgorin functional of the surrounding zeros.

\section{The local collision model}
Place the pair at $x_*\pm h$ and factor
\begin{equation}\label{eq:factor}
  H_0(x)=\bigl[(x-x_*)^2-h^2\bigr]\,q(x),\qquad q(x_*)\neq0,
\end{equation}
where $q$ carries the remaining zeros $\{z_k\}$. Define the inverse distances and
their power sums
\[
  w_k:=\frac{1}{x_*-z_k},\qquad p_1=\sum_k w_k,\quad p_2=\sum_k w_k^2,\quad
  e_2=\tfrac12\bigl(p_1^2-p_2\bigr).
\]

\begin{lemma}[Derivatives at the collision]\label{lem:loc}
With $f(x)=(x-x_*)^2-h^2$ (so $f(x_*)=-h^2$, $f'(x_*)=0$, $f''\equiv2$,
$f'''=f^{(4)}=0$),
\[
  H_0(x_*)=-h^2 q(x_*),\quad
  H_0''(x_*)=2q(x_*)-h^2 q''(x_*),\quad
  H_0^{(4)}(x_*)=12\,q''(x_*)-h^2 q^{(4)}(x_*),
\]
and $q''(x_*)/q(x_*)=p_1^2-p_2=2e_2$.
\end{lemma}
\begin{proof}
Leibniz on $H_0=fq$ with $f''\equiv2$, $f'''=f^{(4)}=0$, $f(x_*)=-h^2$,
$f'(x_*)=0$: the surviving terms are $H_0^{(n)}(x_*)=fq^{(n)}+\binom n2 f''q^{(n-2)}
=\binom n2\!\cdot\!2\,q^{(n-2)}(x_*)-h^2q^{(n)}(x_*)$, giving the three values.
For the last identity, $q'/q=\sum_k w_k=p_1$ and $(q'/q)'=-\sum_k w_k^2=-p_2$, so
$q''/q=(q'/q)'+(q'/q)^2=p_1^2-p_2=2e_2$.
\end{proof}

\begin{proposition}[Gershgorin lift-off law]\label{prop:law}
For an isolated close pair ($h$ small),
\begin{equation}\label{eq:law}
  t_{\mathrm{drown}}=-\frac{H_0(x_*)}{H_0^{(4)}(x_*)}
   =\frac{h^2\,q(x_*)}{12\,q''(x_*)}\bigl(1+O(h^2)\bigr)
   =\frac{h^2}{24\,e_2}\bigl(1+O(h^2)\bigr),
\end{equation}
and the pair lifts off ($t_{\mathrm{drown}}>0$) iff $e_2>0$, i.e.\ iff
$p_1^2>p_2$.
\end{proposition}
\begin{proof}
By Lemma~\ref{lem:loc}, $-H_0(x_*)/H_0^{(4)}(x_*)=h^2q/(12q'')$, and
$q/q''=1/(2e_2)$.
\end{proof}

The lift-off time is the squared half-gap over the second-order crowding $e_2$ of
the surrounding zeros --- a Gershgorin functional. Denser, more one-sided
surroundings (large $e_2$) lift the pair faster.

\section{Why the two criteria coincide}
The sign test $e_2>0\iff p_1^2>p_2$ separates two regimes.
\begin{itemize}
\item \textbf{Symmetric field} (a pure local oscillation): the surrounding zeros
cancel, $p_1=0$, so $e_2=-p_2/2<0$ and $t_{\mathrm{drown}}<0$ --- \emph{no
lift-off}; the floor pushes the extremum further from the axis. (This is the
identity $H_0^{(4)}=\omega^4H_0$ for a pure cosine.) A regular lattice of zeros is
stable under the quartic flow.
\item \textbf{Drift-dominated field} ($\abs{p_1}\gg p_2^{1/2}$): then
$e_2\approx p_1^2/2$ and
\begin{equation}\label{eq:reconcile}
   t_{\mathrm{drown}}\approx\frac{h^2}{12\,p_1^{2}},
\end{equation}
which is the drift law \eqref{eq:drift} with $v\asymp p_1$, the net second-order
crowding $q'/q$ at the site. So the floor and drift criteria are the same law.
\end{itemize}

\begin{remark}
The floor criterion carries the full $e_2=\tfrac12(p_1^2-p_2)$ and is therefore
sharper near the symmetric regime, where $p_2$ is not negligible and the drift
form \eqref{eq:drift} overestimates the lift-off speed. Away from symmetry the two
agree up to a constant.
\end{remark}

\section{Gap compression and a never-emptying band}
The law \eqref{eq:law} converts a statement about gaps into a statement about
times. Suppose a family of configurations has local gaps $h\to0$ while the local
crowding stays bounded, $e_2\asymp c>0$ (a net drift that does not vanish). Then
$t_{\mathrm{drown}}\asymp h^2\to0$: such pairs lift off arbitrarily early.

\begin{conjecture}[Gap-compression instability]\label{conj:gc}
If a real-rooted $H_0$ has, at every sufficiently large scale, pairs with gap
$h\to0$ and bounded-below drift crowding $e_2$, then for every $t>0$ some pair has
already lifted off, so $H_t$ has nonreal zeros for all $t>0$; the real-rooted
region cannot fill the whole line at any finite time.
\end{conjecture}

\noindent This is the natural mechanism by which a fourth-order flow can fail to
restore global real-rootedness: wherever gaps compress faster than any rising
real-rooted window, lift-off descends to meet them. The rigorous inputs it needs
are (i) an effective lower edge for whatever real-rooted window the flow does
produce, and (ii) a lower bound $e_2\ge c>0$ for infinitely many compressing
pairs. Proposition~\ref{prop:law} supplies the explicit rate in between.

\section{The random-matrix spectral-edge instance}
At the spectral edge of a random matrix the surrounding field of eigenvalues is
\emph{one-sided}: all $w_k=1/(x_*-z_k)$ share a sign, so $p_1=\pm\sum_k\abs{w_k}
=\pm S$ with $S=\sum_k\abs{x_*-z_k}^{-1}$, and the cross term cannot cancel.
Hence $e_2=\tfrac12(S^2-p_2)>0$ automatically: \emph{edge pairs are always
unstable inward} under the quartic flow. When the inverse-gaps are spread (no
single neighbor dominates), $p_2/S^2\to0$, so $e_2\approx S^2/2$ saturates its
Maclaurin bound and
\begin{equation}\label{eq:edge}
   t_{\mathrm{drown}}\approx\frac{h^2}{12\,S^2}.
\end{equation}
For a Wigner edge the nearest gap scales as $h\asymp N^{-2/3}$ and $S\asymp N$, so
edge pairs lift off on a time scale $t_{\mathrm{drown}}\asymp N^{-4/3}/N^{2}
=N^{-10/3}$ --- fast, and distributed through the soft-edge gap law: if the
rescaled edge gap has small-gap density $\sim h^{\beta}$ (level repulsion index
$\beta$), the lift-off time inherits the same $\beta$-dependent law. Thus a single
identity, $t_{\mathrm{drown}}\asymp h^2/(\text{crowding})$, governs both
compression settings, with the only difference the gap law that feeds it.

\section*{Scope}
Proposition~\ref{prop:law} is an exact local identity and the reconciliation of
the two criteria is rigorous up to the stated constants. Conjecture~\ref{conj:gc}
and the edge estimate \eqref{eq:edge} are heuristics with explicit rates, not
theorems. We make no claim about the Riemann hypothesis.
\end{document}
